% Unit 2: First Order Linear Differential Equations
% Cheat Sheet (Part I only)
% Prepared by Staples Education
% Content informed by the local curriculum modules (2.1 to 2.5) and the
% QUIZ_DATA micro-practice and unit-mastery items for Unit 2.

\documentclass[11pt]{article}

\usepackage[margin=1in]{geometry}
\usepackage{amsmath}
\usepackage{amssymb}
\usepackage{xcolor}
\usepackage[most]{tcolorbox}
\usepackage{enumitem}
\usepackage{hyperref}
\hypersetup{colorlinks=true, linkcolor=blue!60!black, urlcolor=blue!60!black}

\tcbuselibrary{skins, breakable}

% ---------------------------------------------------------------------------
% Custom colored boxes
% ---------------------------------------------------------------------------

% Blue overview box for the topics summary.
\newtcolorbox{topicsbox}{
    enhanced, breakable,
    colback=blue!5!white, colframe=blue!70!black,
    coltitle=white, fonttitle=\bfseries,
    title=Unit 2 at a Glance: Core Sub-Topics,
    boxrule=0.6pt, arc=2mm,
    left=3mm, right=3mm, top=2mm, bottom=2mm
}

% Orange box for standard forms and master formulas.
\newtcolorbox{formulabox}[1][Master Formula]{
    enhanced, breakable,
    colback=orange!8!white, colframe=orange!80!black,
    coltitle=white, fonttitle=\bfseries,
    title={#1},
    boxrule=0.6pt, arc=1mm,
    left=3mm, right=3mm, top=2mm, bottom=2mm
}

% Purple box with a thick left border for step-by-step worked examples.
\newtcolorbox{examplebox}[1][Worked Example]{
    enhanced, breakable,
    colback=purple!4!white, colframe=purple!70!black,
    coltitle=white, fonttitle=\bfseries,
    title={#1},
    boxrule=0.4pt, sharp corners,
    borderline west={3pt}{0pt}{purple!70!black},
    left=5mm, right=3mm, top=2mm, bottom=2mm
}

\setlist[enumerate]{itemsep=4pt, topsep=4pt}

\begin{document}

% ---------------------------------------------------------------------------
% Title block
% ---------------------------------------------------------------------------
\begin{center}
    {\LARGE \bfseries Unit 2: First Order Linear Differential Equations \\--- Cheat Sheet}\\[6pt]
    {\large Prepared by Staples Education}\\[4pt]
    {\itshape Condensed Cheat Sheet of Core Concepts and Master Formulas}
\end{center}

\vspace{0.4em}

% ---------------------------------------------------------------------------
% Topics overview box
% ---------------------------------------------------------------------------
\begin{topicsbox}
This unit assembles the core solving toolkit of first order theory. Each
sub-topic below builds on the last, moving from raw integration techniques
toward genuine physical models.
\begin{itemize}[leftmargin=1.4em, itemsep=3pt]
    \item \textbf{Separable Equations} --- splitting a rate law into a pure
          $x$ side and a pure $y$ side, then integrating each independently.
    \item \textbf{First-Order Linear Equations and Integrating Factors} ---
          the standard form $y' + P(x)y = Q(x)$ and the factor that collapses
          the left side into a single derivative.
    \item \textbf{Initial Value Problems} --- using a known data point to pin
          the arbitrary constant and select one curve from the solution family.
    \item \textbf{Homogeneous Equations and Substitution} --- the
          substitution $y = vx$ that reduces a degree-matched equation to a
          separable one.
    \item \textbf{Applications: Growth, Decay, and Newton's Law of Cooling}
          --- the exponential models that these methods unlock.
\end{itemize}
\end{topicsbox}

\vspace{0.6em}
\hrule
\vspace{0.6em}

% ===========================================================================
% PART I: CONDENSED CHEAT SHEET
% ===========================================================================
\section*{Part I --- Condensed Cheat Sheet}

\subsection*{2.1 Separable Equations}

A first order equation is \emph{separable} when its right side factors into a
function of $x$ alone times a function of $y$ alone --- that single structural
test is the whole game. Once it factors, you divide and multiply until every
$y$ rides with $dy$ and every $x$ rides with $dx$, then integrate both sides.

\begin{formulabox}[Standard Form and Method: Separable]
A separable equation and its solution scheme:
\[
    \frac{dy}{dx} = g(x)\,h(y)
    \quad\Longrightarrow\quad
    \int \frac{1}{h(y)}\,dy = \int g(x)\,dx + C.
\]
One constant $C$ suffices, since the two constants of integration merge into
a single arbitrary constant.
\end{formulabox}

\noindent\textbf{Pitfall.} Dividing by $h(y)$ silently assumes $h(y) \neq 0$.
Any constant value of $y$ that makes $h(y) = 0$ is an equilibrium solution that
this step can quietly discard --- always check it separately.

\begin{examplebox}[Worked Example 2.1: Solve $\frac{dy}{dx} = 6x^{2}y$]
\textbf{Step 1.} Separate the variables:
\[ \frac{1}{y}\,dy = 6x^{2}\,dx. \]
\textbf{Step 2.} Integrate both sides:
\[ \ln|y| = 2x^{3} + C. \]
\textbf{Step 3.} Exponentiate to isolate $y$:
\[ y = e^{2x^{3} + C} = C\,e^{2x^{3}}. \]
The dropped solution $y = 0$ is recovered as the case $C = 0$.
\end{examplebox}

\subsection*{2.2 First-Order Linear Equations and Integrating Factors}

A first order equation is \emph{linear} when $y$ and $y'$ each appear only to
the first power and never multiply one another. Before reading off anything,
you must normalize so the derivative carries coefficient one --- only then is
$P(x)$ truly the coefficient of $y$.

\begin{formulabox}[Standard Form and Integrating Factor]
Standard form, integrating factor, and solution:
\[
    \frac{dy}{dx} + P(x)\,y = Q(x),
    \qquad
    \mu(x) = e^{\int P(x)\,dx},
\]
\[
    \frac{d}{dx}\big[\mu(x)\,y\big] = \mu(x)\,Q(x)
    \quad\Longrightarrow\quad
    y = \frac{1}{\mu(x)}\left( \int \mu(x)\,Q(x)\,dx + C \right).
\]
\end{formulabox}

\noindent The magic of $\mu(x)$ is that multiplying through turns the entire
left side into the exact product-rule derivative $\frac{d}{dx}[\mu y]$ ---
that is the single fact the whole method rests on.

\begin{examplebox}[Worked Example 2.2: Solve $\frac{dy}{dx} + y = e^{x}$]
\textbf{Step 1.} Read off $P(x) = 1$, so the integrating factor is
\[ \mu = e^{\int 1\,dx} = e^{x}. \]
\textbf{Step 2.} Multiply through; the left side collapses to a derivative:
\[ \frac{d}{dx}\big[e^{x}y\big] = e^{x}\cdot e^{x} = e^{2x}. \]
\textbf{Step 3.} Integrate both sides:
\[ e^{x}y = \frac{1}{2}e^{2x} + C. \]
\textbf{Step 4.} Divide by $e^{x}$:
\[ y = \frac{1}{2}e^{x} + C e^{-x}. \]
\end{examplebox}

\subsection*{2.3 Initial Value Problems}

The general solution of a first order equation is a whole family of curves
indexed by one arbitrary constant. An \emph{initial condition} $y(x_0) = y_0$
selects exactly one member --- it fixes the constant to the value that forces
the curve through the given point.

\begin{formulabox}[General versus Particular Solution]
A first order equation needs exactly \textbf{one} initial condition. The
general solution carries the free constant; the particular solution has it
determined:
\[
    y = F(x) + C \;\;\xrightarrow{\;y(x_0)=y_0\;}\;\; C = y_0 - F(x_0).
\]
\end{formulabox}

\begin{examplebox}[Worked Example 2.3: Solve $\frac{dy}{dx} = 3x^{2}$, $\;y(1)=5$]
\textbf{Step 1.} Integrate for the general solution:
\[ y = x^{3} + C. \]
\textbf{Step 2.} Apply $y(1) = 5$:
\[ 5 = (1)^{3} + C \;\Longrightarrow\; C = 4. \]
\textbf{Step 3.} State the particular solution:
\[ y = x^{3} + 4. \]
\end{examplebox}

\subsection*{2.4 Homogeneous First-Order Equations and Substitution}

A function $f(x,y)$ is \emph{homogeneous of degree $n$} when scaling both
inputs pulls out a clean power of the scale factor: $f(tx, ty) = t^{n}f(x,y)$.
An equation $M\,dx + N\,dy = 0$ is homogeneous when $M$ and $N$ share the same
degree --- and that shared degree is precisely what lets the ratio
substitution work.

\begin{formulabox}[The Substitution $y = vx$]
For a homogeneous equation, set
\[
    y = vx, \qquad \frac{dy}{dx} = v + x\frac{dv}{dx},
\]
which always reduces the equation to a \emph{separable} one in $v$ and $x$.
After solving for $v$, back-substitute $v = \dfrac{y}{x}$.
\end{formulabox}

\begin{examplebox}[Worked Example 2.4: Solve $\frac{dy}{dx} = \frac{x + 2y}{x}$]
\textbf{Step 1.} Rewrite the right side and substitute $y = vx$:
\[ \frac{dy}{dx} = 1 + 2\frac{y}{x} = 1 + 2v, \qquad v + x\frac{dv}{dx} = 1 + 2v. \]
\textbf{Step 2.} Cancel $v$ and separate:
\[ x\frac{dv}{dx} = 1 + v \;\Longrightarrow\; \frac{1}{1+v}\,dv = \frac{1}{x}\,dx. \]
\textbf{Step 3.} Integrate and exponentiate:
\[ \ln|1+v| = \ln|x| + C \;\Longrightarrow\; 1 + v = Cx. \]
\textbf{Step 4.} Back-substitute $v = \frac{y}{x}$ and clear:
\[ 1 + \frac{y}{x} = Cx \;\Longrightarrow\; y = Cx^{2} - x. \]
\end{examplebox}

\subsection*{2.5 Applications: Growth, Decay, and Newton's Law of Cooling}

When a rate of change is proportional to the current amount, the solution is
exponential --- this single idea drives population models, radioactive decay,
and cooling alike.

\begin{formulabox}[Exponential Models]
Growth and decay:
\[
    \frac{dP}{dt} = kP \;\Longrightarrow\; P(t) = P_{0}e^{kt}
    \qquad (k>0 \text{ grows},\; k<0 \text{ decays}),
\]
with doubling time $t = \dfrac{\ln 2}{k}$. Newton's law of cooling:
\[
    \frac{dT}{dt} = -k\,(T - T_{s}) \;\Longrightarrow\;
    T(t) = T_{s} + (T_{0} - T_{s})\,e^{-kt},
\]
so $T \to T_{s}$ as $t \to \infty$.
\end{formulabox}

\begin{examplebox}[Worked Example 2.5: A culture doubles in 3 hours]
\textbf{Step 1.} The model $P = P_{0}e^{kt}$ doubles when $e^{3k} = 2$:
\[ 3k = \ln 2 \;\Longrightarrow\; k = \frac{\ln 2}{3}. \]
\textbf{Step 2.} The growth law is therefore
\[ P(t) = P_{0}\,e^{(\ln 2 / 3)\,t} = P_{0}\,2^{\,t/3}. \]
Every $3$ hours the population multiplies by two, as required.
\end{examplebox}

\vspace{0.8em}
\begin{center}
\textit{End of Unit 2 Cheat Sheet. Prepared by Staples Education.}
\end{center}

\end{document}
